\chapter{1954年全国卷}

\section{解答题}

\begin{problem}
\begin{enumerate}
\item 化简 \(\left[\left(a^{-\frac{3}{2}}\cdot b^{2}\right)^{-1}\cdot\left(a\cdot b^{-3}\right)^{\frac{1}{2}}\cdot\left(b^{\frac{1}{2}}\right)^{7}\right]^{\frac{1}{3}}\).
\item 解 \(\frac{1}{6}\log x=\frac{1}{3}\log a+2\log b+\log c\).
\item 用二项式定理计算 \(\left(3.02\right)^{4}\)，使误差小于 \(\frac{1}{1000}\).
\item 直角三角形弦上半圆的面积等于勾上半圆与股上半圆面积之和，试证明之.
\item 已知球的半径为 \(r\)，求内接正方体的体积.
\item 已知三角形的一边之长为 \(a\)，两邻角为 \(\beta\) 及 \(\gamma\)，求计算边长 \(b\) 的计算公式.
\end{enumerate}
\end{problem}

\begin{solution}
\begin{enumerate}
\item 在实数范围内原式有意义时 \(a>0\)，\(b>0\). 因此
\[
\left[\left(a^{-\frac{3}{2}}b^{2}\right)^{-1}\left(ab^{-3}\right)^{\frac{1}{2}}\left(b^{\frac{1}{2}}\right)^{7}\right]^{\frac{1}{3}}
=\left(a^{\frac{3}{2}}b^{-2}a^{\frac{1}{2}}b^{-\frac{3}{2}}b^{\frac{7}{2}}\right)^{\frac{1}{3}}
=\left(a^{2}b^{0}\right)^{\frac{1}{3}}=a^{\frac{2}{3}}
\]

\item 对数的真数必须为正，故 \(x>0\)，\(a>0\)，\(b>0\)，\(c>0\). 由对数的运算性质，原方程等价于
\[
\log\left(x^{\frac{1}{6}}\right)=\log\left(a^{\frac{1}{3}}b^{2}c\right)
\]
由于对数函数在其定义域上具有单调性，得 \(x^{\frac{1}{6}}=a^{\frac{1}{3}}b^{2}c\). 两边均为正数，六次方后得到 \(x=\left(a^{\frac{1}{3}}b^{2}c\right)^{6}=a^{2}b^{12}c^{6}\).

\item 由二项式定理，
\[
\left(3.02\right)^{4}=\left(3+0.02\right)^{4}
=3^{4}+4\cdot3^{3}\cdot0.02+6\cdot3^{2}\cdot\left(0.02\right)^{2}
+4\cdot3\cdot\left(0.02\right)^{3}+\left(0.02\right)^{4}
\]
从第 4 项起的余项为 \(4\cdot3\cdot\left(0.02\right)^{3}+\left(0.02\right)^{4}=0.00009616<0.001\)，所以取前三项所得近似值的误差小于 \(\frac{1}{1000}\). 计算前三项得
\[
3^{4}+4\cdot3^{3}\cdot0.02+6\cdot3^{2}\cdot\left(0.02\right)^{2}=81+2.16+0.0216=83.1816
\]
因此可取 \(\left(3.02\right)^{4}\approx83.182\)，且 \(\left|83.182-\left(3.02\right)^{4}\right|=0.00030384<0.001\).

\item 设直角三角形的两直角边，即勾和股，长分别为 \(a\)，\(b\)，斜边，即弦，长为 \(c\). 由勾股定理，\(c^{2}=a^{2}+b^{2}\). 直径为 \(d\) 的半圆面积为 \(\frac{\pi d^{2}}{8}\)，故
\[
\frac{\pi c^{2}}{8}=\frac{\pi\left(a^{2}+b^{2}\right)}{8}=\frac{\pi a^{2}}{8}+\frac{\pi b^{2}}{8}
\]
左边是弦上半圆的面积，右边两项分别是勾上半圆和股上半圆的面积，所以命题得证.

\item 设内接正方体的棱长为 \(s\). 正方体的外接球球心是正方体中心，正方体的空间对角线两端是球面上的顶点，因此空间对角线等于球的直径 \(2r\). 由正方体空间对角线公式，\(\sqrt{3}s=2r\)，从而 \(s=\frac{2r}{\sqrt{3}}=\frac{2\sqrt{3}}{3}r\). 所以内接正方体的体积为
\[
V=s^{3}=\left(\frac{2r}{\sqrt{3}}\right)^{3}=\frac{8\sqrt{3}}{9}r^{3}
\]

\item 三角形内角满足 \(\beta>0\)，\(\gamma>0\)，\(\beta+\gamma<180^{\circ}\). 已知边 \(a\) 的对角为 \(180^{\circ}-\left(\beta+\gamma\right)\)，边 \(b\) 的对角为 \(\beta\). 由正弦定理，
\[
\frac{a}{\sin\left[180^{\circ}-\left(\beta+\gamma\right)\right]}=\frac{b}{\sin\beta}
\]
所以
\[
b=\frac{a\sin\beta}{\sin\left[180^{\circ}-\left(\beta+\gamma\right)\right]}=\frac{a\sin\beta}{\sin\left(\beta+\gamma\right)}
\]
\end{enumerate}
\end{solution}

\begin{problem}
描绘 \(y=3x^{2}-7x-1\) 之图象，并按下列条件分别求 \(x\) 的值的范围：

\begin{enumerate}
\item \(y>0\)；
\item \(y<0\).
\end{enumerate}
\end{problem}

\begin{solution}
配方得
\[
y=3\left(x-\frac{7}{6}\right)^{2}-\frac{61}{12}
\]
即 \(\left(x-\frac{7}{6}\right)^{2}=\frac{1}{3}\left(y+\frac{61}{12}\right)\). 因此抛物线开口向上，对称轴为 \(x=\frac{7}{6}\)，顶点为 \(V\left(\frac{7}{6},-\frac{61}{12}\right)\)，且与 \(y\) 轴交于 \((0,-1)\). 令 \(y=0\)，得
\(3x^{2}-7x-1=0\)，\(x=\frac{7\pm\sqrt{61}}{6}\).
所以抛物线与 \(x\) 轴的交点为 \(M\left(\frac{7-\sqrt{61}}{6},0\right)\) 和 \(N\left(\frac{7+\sqrt{61}}{6},0\right)\).

\begin{enumerate}
\item 因为抛物线开口向上，故 \(y>0\) 时 \(x\) 在两个零点之外，即 \(x\in\left(-\infty,\frac{7-\sqrt{61}}{6}\right)\cup\left(\frac{7+\sqrt{61}}{6},+\infty\right)\).

\item 因为二次项系数 \(3>0\)，所以 \(y<0\) 时 \(x\) 在两个零点之间，即 \(x\in\left(\frac{7-\sqrt{61}}{6},\frac{7+\sqrt{61}}{6}\right)\).
\end{enumerate}

\solutionfigure{\bitmapfigure[max width=0.55\linewidth]{img/1954/national/q2_solution_fig1.png}}
\end{solution}

\begin{problem}
假设两圆互相外切，求证用连心线段为直径所作的圆必与前两圆的外公切线相切.
\end{problem}

\begin{solution}
设 \(\odot O_{1}\) 和 \(\odot O_{2}\) 的半径分别为 \(r_{1}\) 和 \(r_{2}\)，两圆外切. 设一条外公切线为 \(A_{1}A_{2}\)，切点分别为 \(A_{1}\) 和 \(A_{2}\). 令 \(O\) 为 \(O_{1}O_{2}\) 的中点，并过 \(O\) 作 \(OA\perp A_{1}A_{2}\)，其中 \(A\) 是垂足. 因为外公切线与圆的半径垂直，且两个圆心在该切线的同侧，所以中点到该直线的距离为两圆心到该直线距离的平均值，从而
\[
OA=\frac{O_{1}A_{1}+O_{2}A_{2}}{2}=\frac{r_{1}+r_{2}}{2}
\]
两圆外切时，圆心距等于两半径之和，即 \(O_{1}O_{2}=r_{1}+r_{2}\)，故 \(OA=\frac{1}{2}O_{1}O_{2}\). 以 \(O_{1}O_{2}\) 为直径所作的圆，其圆心为 \(O\)，半径为 \(\frac{1}{2}O_{1}O_{2}=OA\). 由于圆心到直线 \(A_{1}A_{2}\) 的距离等于该圆半径，直线 \(A_{1}A_{2}\) 与此圆相切. 对另一条外公切线作同样论证，即得此圆也与另一条外公切线相切.

\solutionfigure{\bitmapfigure[max width=0.55\linewidth]{img/1954/national/q3_solution_fig1.png}}
\end{solution}

\begin{problem}
解 \(\frac{1+\tan x}{1-\tan x}=1+\sin 2x\)，求 \(x\) 之通值.
\end{problem}

\begin{solution}
原方程有意义的条件是 \(\cos x\ne0\) 且 \(1-\tan x\ne0\)，也就是 \(\cos x\ne0\) 且 \(\cos x-\sin x\ne0\). 在此条件下，利用 \(\sin 2x=2\sin x\cos x\)，有
\(\frac{1+\tan x}{1-\tan x}=\frac{\cos x+\sin x}{\cos x-\sin x}\)，\(1+\sin 2x=\left(\cos x+\sin x\right)^{2}\).
因此
\[
\cos x+\sin x=\left(\cos x+\sin x\right)^{2}\left(\cos x-\sin x\right)
\]
即
\[
\left(\cos x+\sin x\right)\left[1-\left(\cos x+\sin x\right)\left(\cos x-\sin x\right)\right]=0
\]
由于 \(\left(\cos x+\sin x\right)\left(\cos x-\sin x\right)=\cos^{2}x-\sin^{2}x\)，上式化为
\[
2\left(\cos x+\sin x\right)\sin^{2}x=0
\]
于是分两类求解：\(\cos x+\sin x=0\) 时，\(\tan x=-1\)，所以 \(x=k\pi-\frac{\pi}{4}\)；\(\sin^{2}x=0\) 时，\(x=k\pi\)，其中 \(k\in\Z\). 对这两类值分别有 \(\tan x=-1\) 或 \(\tan x=0\)，均满足原方程的定义域，代入原方程也分别得到 \(0=0\) 和 \(1=1\). 故方程的通解为 \(x=k\pi-\frac{\pi}{4}\) 或 \(x=k\pi\)，其中 \(k\in\Z\).
\end{solution}

\begin{problem}
有一直圆锥，全面积为 \(a\)；与之同底同高之直圆柱全面积为 \(a'\). 求该圆锥高与母线之比.
\end{problem}

\begin{solution}
设直圆锥的高为 \(h\)，底面半径为 \(R\)，母线长为 \(l\)，则 \(l^{2}=R^{2}+h^{2}\)，且 \(0<h<l\). 圆锥和同底同高圆柱的全面积分别为
\(a=\pi R\left(R+l\right)\)，\(a'=2\pi R\left(R+h\right)\).
因此
\[
2a\left(R+h\right)=a'\left(R+l\right)
\]
令 \(u=\frac{h}{l}\)，并记 \(D=2a-a'\). 由 \(R=\sqrt{l^{2}-h^{2}}\) 和上式，得
\[
D\sqrt{1-u^{2}}=a'-2au
\]
又由几何关系，
\[
D=2a-a'=2\pi R(l-h)>0
\]
此外，
\[
a'-a=\pi R\left(R+2h-l\right)>0
\]
因为 \((R+2h)^{2}-l^{2}=4Rh+3h^{2}>0\)，所以 \(R+2h-l>0\). 因而 \(a<a'<2a\)，从而 \(0<D<a\). 原方程左边为正，故右边也必须为正，即 \(a'-2au>0\). 两边平方，得到
\[
\left(4a^{2}+D^{2}\right)u^{2}-4aa'u+4a\left(a'-a\right)=0
\]
其判别式为
\[
\Delta=\left(-4aa'\right)^{2}-4\left(4a^{2}+D^{2}\right)4a\left(a'-a\right)=16aD^{3}
\]
所以平方后所得两根为
\[
u_{\pm}=\frac{2aa'\pm2D\sqrt{aD}}{4a^{2}+D^{2}}
\]
平方可能引入增根，必须回代未平方的方程. 对上述两根，
\[
a'-2au_{\pm}=\frac{D\left(Da'\mp4a\sqrt{aD}\right)}{4a^{2}+D^{2}}
\]
由于 \(0<D<a\) 且 \(a'<2a\)，有
\[
D(a')^{2}<a(2a)^{2}<16a^{3}
\]
即 \(Da'<4a\sqrt{aD}\). 因此 \(u_{+}\) 使 \(a'-2au_{+}<0\)，不满足未平方的方程.

另一方面，\(a'>a>D>0\)，从而 \(a'>\sqrt{aD}\) 且 \(a>D\)，所以 \(aa'>D\sqrt{aD}\)，即 \(u_{-}>0\)，且
\(u_{-}<\frac{2aa'}{4a^{2}+D^{2}}<1\)，\(4a^{2}+D^{2}-2aa'=D(2a+D)>0\).
对 \(u_{-}\)，上式中的 \(a'-2au_{-}\) 为正；由于它满足平方后的方程，故同时满足原来的未平方方程. 所以唯一可行的比值为
\[
\boxed{\frac{h}{l}=\frac{2aa'-2\left(2a-a'\right)\sqrt{a\left(2a-a'\right)}}{4a^{2}+\left(2a-a'\right)^{2}}}
\]
\end{solution}
